\section{Evaluation Benchmark}
\label{sec:benchmark}



We establish a comprehensive evaluation benchmark covering 
camera pose estimation and 3D reconstruction across diverse indoor, 
outdoor, and large-scale environments.

\subsection{Datasets}
\label{sec:eval_datasets}

Streaming 3D reconstruction must generalize across a wide spectrum of
scenarios, from object-centric captures to room-scale interiors and
city-scale outdoor trajectories, under varying sequence lengths,
camera motions, and scene complexities.
To thoroughly evaluate this, our benchmark is built on five
complementary datasets: Oxford Spires~\cite{tao2025spires},
ETH3D~\cite{schops2017multi}, 7-Scenes~\cite{shotton2013scene},
Tanks and Temples~\cite{Knapitsch2017}, and NRGBD~\cite{azinovic2022neural},
which collectively cover indoor and outdoor environments,
object-centric and multi-scene trajectories, short sequences
(tens of frames) to long sequences (thousands of frames), and
both synthetic-grade and real-world capture conditions.
Below we describe each dataset and our evaluation configuration.

\noindent\textbf{Oxford Spires}~\cite{tao2025spires} is a large-scale
outdoor and indoor dataset captured across the historic Oxford campus,
featuring complex scene transitions, revisits, and significant scale
variation within each sequence.
The dataset provides ground-truth camera trajectories from a
high-precision LiDAR-inertial SLAM system.
We select 13 scenes with available ground-truth trajectories: Keble
College~02--05, Observatory Quarter~01--02, Blenheim Palace~01--02
and~05, Christ Church~01--03 and~05, and Bodleian Library~02.
We use camera~0 as the viewpoint and evaluate under two settings:
\textit{sparse} (320 frames, sampled every 12 frames) to test
single-pass reconstruction within our training range, and
\textit{dense} (3{,}840 frames) to stress-test long-sequence
streaming capabilities.

\noindent\textbf{ETH3D}~\cite{schops2017multi} provides
high-resolution indoor and outdoor images with ground-truth depth maps
acquired from a laser scanner.
The scenes span offices, lecture rooms, relief structures, facades,
playgrounds, terraces, and botanical gardens, covering both small-scale
indoor environments and large-scale outdoor settings.
We follow the evaluation configurations established in
DA3~\cite{depthanything3}, using all available frames.
We use a threshold of $d = 0.1$\,m for the F1 reconstruction metric.

\noindent\textbf{7-Scenes}~\cite{shotton2013scene} is a widely-used
indoor RGB-D dataset consisting of 7 scenes (Chess, Fire, Heads,
Office, Pumpkin, Kitchen, Stairs) captured with a Kinect sensor at
$640 \times 480$ resolution.
It is a challenging real-world dataset: the images contain significant
motion blur, and the scenes feature textureless surfaces and repetitive
structures that are difficult for pose estimation.
We downsample the number of frames for each scene by a stride of 5
to reduce redundancy from the high-framerate Kinect capture while
retaining sufficient viewpoint coverage.

\noindent\textbf{Tanks and Temples}~\cite{Knapitsch2017} is a
large-scale outdoor dataset providing high-resolution images, depth
maps from LiDAR, and ground-truth 3D shapes for benchmarking
multi-view reconstruction.
We select 6 scenes covering diverse structures: Barn, Caterpillar,
Church, Ignatius, Meeting Room, and Truck, with all images included.

\noindent\textbf{NRGBD}~\cite{azinovic2022neural} is a dataset
designed for neural RGB-D surface reconstruction, containing indoor
scenes with high-quality ground-truth depth from structured-light
sensors.
The scenes include cluttered room-scale environments with fine
geometric details.
We sample images with a stride of 5 and evaluate dense reconstruction
quality using the F1 metric.

\subsection{Metrics}
\label{sec:metrics}

\noindent\textbf{Camera Pose Estimation.}
We evaluate camera pose accuracy using the Area Under the Curve (AUC) of the relative pose error at angular thresholds of $3\degree$ and $30\degree$.
For trajectory-level evaluation, we report the Absolute Trajectory Error (ATE) in meters, which measures global trajectory consistency after Sim(3) alignment, as well as Relative Pose Error for translation (RPE-trans) and rotation (RPE-rot), which capture local frame-to-frame accuracy.

\noindent\textbf{3D Reconstruction.} 
We evaluate reconstruction quality using the F1 score, Accuracy (Acc), and Completeness (Comp).
The reconstructed and ground-truth point clouds are first aligned using the Umeyama~
\cite{umeyama2002least} method, followed by ICP refinement (threshold 0.1).
For ETH3D, which features large-scale scenes with relatively sparse ground-truth points, we adopt an F1 threshold of 0.25 with a voxel size of 0.039m.
For 7-Scenes and NRGBD, which contain many frames and dense depth projections producing tens of millions of points, we first downsample both point clouds to a uniform voxel grid (voxel size $4.0/512$) before ICP, and set the F1 threshold to 0.05.



